\subsection{Four named obstructions}
\label{subsec:four-named-obstructions-ch5}

The eighth entry of the Dirac--Igusa pro-object is the obstruction
theory.  The cross-level identity ②$\leftrightarrow$③ between the
imported Borcherds--Kac--Moody denominator algebra and the compact
protected Pfaffian passes through four obstruction classes and the
finite geometric Pfaffian datum.  The first obstruction is reduced on
$K3\times E$ to the retained
D0-HN datum in
Appendix~\ref{app:obstruction-discharges}
(Theorem~\ref{thm:G-D0}); \((O1)\) and \((O1)^+\) are reduced there
to the retained quotient-orientation and Weyl-transport data
(Theorems~\ref{thm:G-O1}, \ref{thm:G-O1-plus}); the fourth is reduced
to the retained type-II
wall-atlas datum \((O2_{\mathrm{atlas}})\) by
Theorem~\ref{thm:G-O2-orbit-transport}.  The finite geometric Pfaffian
datum \((P_{\mathrm{fin}})\) of
Definition~\ref{def:finite-geometric-pfaffian-datum} supplies the
cosection-reduced skew perfect complexes, Pfaffian line, finite
Pfaffian section, type-II wall divisor order, automorphic
line-comparison, support and parity rows, leading coefficient, and
Pfaffian-line Mittag--Leffler compatibility for the cofinal HN product.
The principal
Pfaffian and orientation theorem
\[
\operatorname{Pf}_{\mathrm{prot}}(\mathfrak D_X^{\mathrm{DI}})
=\Delta_5,
\qquad
\varepsilon_o=\nu_{\Delta_5}
\]
holds relative to the first five components of the retained
obstruction datum
\(\mathfrak R_X^{\mathrm{DI}}\), namely
\((D0_{\mathrm{HN}},O1_{\mathrm{quot}},\mathcal T_W,
O2_{\mathrm{atlas}},P_{\mathrm{fin}})\).
The Lie-superalgebra recognition \(P_X\cong\mathfrak g_{\Delta_5}\)
requires the finite compact-source recognition datum \((S1)\)--\((S10)\)
of \S\ref{subsec:primitive-recognition}.
On \(K3\times E\) the D0 criterion, the quotient-orientation criteria,
and the conditional wall-atlas transport are stated in
Appendix~\ref{app:obstruction-discharges}; the compact-source
recognition datum is a separate hypothesis.
The retained obstruction data are the components of the tuple
\(\mathfrak R_X^{\mathrm{DI}}\) of
Definition~\ref{def:G-retained-obstruction-datum}.  Thus
``retained datum'' in this section means a projection of that tuple,
not a scalar normalization or a theorem label.

The fourth obstruction is the type-II wall-atlas datum; the
Appendix~G theorem transports that datum along the Weyl orbit and
compares its character with \(\nu_{\Delta_5}\).

\subsubsection*{(D0) D0-degeneration limit of the cosection-reduced
d-critical orientation theory}

\textit{Obstruction.}  At every finite HN height $R$ and every retained
charge stratum $\mathcal C\in\mathcal C_{\sigma,S,\le R}$, the
cosection-reduced d-critical structure on the K3-fibre direction must
extend through a retained flat \(D0\)-degeneration family to the
zero-dimensional sector on \(K3\times E\).  Equivalently, the datum
\((D0_{\mathrm{HN}})\) of
Definition~\ref{def:G-D0-HN-system} supplies the finite derived
degeneration family, the HN transition morphisms, the reduced
obstruction theory, the additive K3 semiregularity cosection
\(\operatorname{CS}_{\mathcal C,R}\), the reduced vanishing-cycle
specialization, the Joyce--Upmeier orientation specialization, the
Pfaffian-line and protected-integration specialization, and strict
Mittag--Leffler exactness.  The filtered Hall additivity equation
\[
q^*\operatorname{CS}_C
=p_1^*\operatorname{CS}_A+p_2^*\operatorname{CS}_B
\]
is part of this datum on extension stacks throughout the limit.
The local/wrapped boundary \(b_R^{\mathrm{geom}}\to0\) belongs to the
hybrid Ran carrier; it enters \((D0)\) only through compatibility of
the retained Hall correspondences with the zero-dimensional limit.

\textit{Finite site.}  The finite Hall atlas and hybrid carrier of
\S\ref{subsec:eight-finite-constructions}.  Failure of (D0) prevents
construction of the reduced Hall product $m_R^{\mathrm{red}}$, the
orientation gerbe descent of (O1), and the Pfaffian line of (D0-Pf)
on the wrapped sector.

\textit{Criterion.}  Criterion on $K3\times E$ in
Theorem~\ref{thm:G-D0}, relative to the retained D0-HN datum of
Definition~\ref{def:G-D0-HN-system}.  The
fibrewise K3 framework of \cite{JoyceSong,JoyceUpmeier} supplies the
orientation technology on the projection-finite stratum; the wrapped
stratum limit is part of the retained datum of
Appendix~\ref{app:G-D0-discharge}.
The finite D0 proof tables of
Definition~\ref{def:G-D0-HN-system} are the checkable form of this
datum.  A Hilbert-scheme scalar specialization, a K3 scalar test, or a
scalar trace row is not a D0-HN transition row.

\subsubsection*{(O1) Strong reduced orientation on the
$K3\times E$-quotient self-Ext frame bundle}

\textit{Obstruction.}  On the cosection-reduced moduli
$\mathfrak M^{\mathrm{red}}_{\mathcal C,R}$ at every finite HN height
$R$, the orientation line $o_{\mathcal C}$ has chosen square root
$o_{\mathcal C}^{\otimes 2}\simeq\det K^{\mathrm{red}}_{\mathcal C}$
\cite{JoyceUpmeier}, multiplicative under reduced extension
correspondences, with the quotient-descent equations: vanishing of the
reduced gerbe class
\[
\alpha^{\mathrm{red}}_{\mathcal C}
=w_2(\det K^{\mathrm{red}}_{\mathcal C})
+\mathrm{cs}^{*}w_2(\det\operatorname{coker}\mathrm{cs})
\in H^2(\mathfrak M^{\mathrm{red}}_{\mathcal C};\mathbb F_2),
\]
the free $E$-class
$\alpha^{E,\mathrm{free}}_{\mathcal C}=
a_1(\mathcal C)u_1+a_2(\mathcal C)u_2\in H^2(BE;\mathbb F_2)$,
and, for each finite-stabilizer stratum $S$ with stabilizer
$H\subset E_{\mathrm{tors}}$, the equivariant gerbe class
$\widetilde\beta^{H}_{\mathcal C,S}\in H^2_H(S;\mathbb F_2)$ together
with the choice of zero $H^1(BH;\mathbb F_2)$-linearization character
$\lambda^H_{\mathcal C,S}=0$.  The statement applies on every charge
stratum used by the Hall correspondences and is compatible with HN
transitions.

\textit{Finite site.}  The orientation descent datum of
\S\ref{subsec:eight-finite-constructions}; (O1) datum of
\S\ref{subsec:cosection-twisted-orientation}.  Failure of (O1)
prevents quotient-after-correspondence descent and obstructs the
reduced Hall product, the Pfaffian line on the reduced quotient,
and the orientation character.

\textit{Criterion.}  Theorem~\ref{thm:G-O1} is a
quotient-orientation criterion relative to the retained
\((O1_{\mathrm{quot}})\) datum.  The Klein-four condition at
$H=E[2]$ reduces the edge class to \(r_1=r_2=r_3=0\), but quotient
orientation also requires the zero degree-one linearization character.
For even $N\ge 4$ the full two-primary residual includes the mixed
coefficient \(A_{12}^{(N)}\), which cyclic restrictions do not detect.

\subsubsection*{(O1)$^+$ Weyl-equivariant transport along
$W^{(2)}(\Lambda_{II}^{2,1})$}

\textit{Obstruction.}  The orientation line of (O1) carries
Weyl-equivariant lifts of the three type-II wall transports
$\tau_i:o_{\mathcal C}\to s_{\delta_i}^* o_{\mathcal C}$
($i=1,2,3$) such that:
\begin{enumerate}
\item the $2$-cocycle $c_{o,R}\in Z^2(\mathcal G_R;
\underline{\mathbb F}_2)$ on the finite partial action groupoid
$\mathcal G_R$ has vanishing cohomology class
$[c_{o,R}]=0$ in
$H^2(W^{(2)}(\Lambda_{II}^{2,1});\mathbb F_2)$ on each complete
orbit, with chosen $1$-cochain killing $c_{o,R}$;
\item $\tau_i^2=1$ for $i=1,2,3$ and the Coxeter relations of
$W^{(2)}(\Lambda_{II}^{2,1})$ hold on the lifts;
\item the quotient-orientation cocycle $\mathfrak q_{\mathcal C}$
transports along $\tau_i$ between charge strata over Weyl-related
Gram degrees, with vanishing torsor defect
$\omega_{i,\mathcal C}=0$.
\end{enumerate}

\textit{Finite site.}  The Weyl-transport datum of
\S\ref{subsec:eight-finite-constructions}; (O1)$^+$ datum of
\S\ref{subsec:cosection-twisted-orientation}.  Failure of (O1)$^+$
obstructs the orientation character $\varepsilon_o$ and prevents
recognition of the Weyl orientation pattern as
$\nu_{\Delta_5}|_{W^{(2)}(\Lambda_{II}^{2,1})}$.

\textit{Conditional theorem.}  The equivariance part of
Theorem~\ref{thm:G-O1-plus} is conditional on the retained
quotient-orientation datum and chosen Weyl lifts.
The local rank-one Pfaffian crossing
supplies the local sign $\varepsilon^{\mathrm{loc}}_{\delta_i,R}=-1$
on a generic chart; the character comparison also uses the retained
\((O2_{\mathrm{atlas}})\) datum.  The global Coxeter coherence on
$W^{(2)}(\Lambda_{II}^{2,1})$ orbits is established by the
no-braid type-II Coxeter presentation and Maass-character match of
Appendix~\ref{app:G-O1-plus-discharge}.

\subsubsection*{(O2) Local Pfaffian wall normal form on type-II
walls}

\textit{Obstruction.}  At every finite HN height $R$ and every retained
component, boundary stratum, and chart overlap of the three simple
type-II walls $\Hpl_{\delta_i}$, $i=1,2,3$, the cosection-reduced
self-Ext complex splits as
\[
K^{\mathrm{red}}_{\delta_i,R}\simeq
K^{\mathrm{tan}}_{\delta_i,R}\oplus K^{\mathrm{nor}}_{\delta_i,R},
\]
with $s_{\delta_i}$-equivariant tangent factor and the rank-one
normal form
\[
K^{\mathrm{nor}}_{\delta_i,R}\simeq
[\R\xrightarrow{u_{\delta_i,R}}\R],
\qquad
N_{\delta_i,R}^{\mathrm{Pf}}=1,
\]
with normal Pfaffian unit
$\operatorname{Pf}(K^{\mathrm{nor}}_{\delta_i,R})
=\upsilon_{\delta_i,R}u_{\delta_i,R}$,
$s_{\delta_i}(\upsilon_{\delta_i,R})=\upsilon_{\delta_i,R}$,
$s_{\delta_i}(u_{\delta_i,R})=-u_{\delta_i,R}$.  The normal Pfaffian
units, tangent/normal splittings, quotient orientations, and
protected integration maps agree on chart overlaps.  The half-Hilbert
orbit sign sum of size $4096=2^{12}$ on the retained components and
boundary strata of the $W^{(2)}(\Lambda_{II}^{2,1})$-orbit reproduces
the reflection character $\nu_{\Delta_5}|_{W^{(2)}(\Lambda_{II}^{2,1})}$.

\textit{Finite site.}  The Pfaffian wall-atlas datum of
\S\ref{subsec:eight-finite-constructions}; rank-one Pfaffian crossing
of \S\ref{subsec:dirac-pfaffian}.  The finite proof surface is
\(\mathrm{wall\_objects}\), \(\mathrm{wall\_charts}\),
\(\mathrm{overlaps}\), \(\mathrm{orbit\_transport}\),
\(\mathrm{half\_hilbert\_orbit}\), \(\mathrm{compatibilities}\), and
\(\mathrm{scalar\_separation}\); these rows distinguish the rank-one
local sign, the global half-Hilbert atlas, and the scalar equality
\(2^{12}=64^2\).  Failure of (O2) prevents the
local Pfaffian section
$\operatorname{pf}_{\delta_i,R}=\upsilon_{\delta_i,R}u_{\delta_i,R}$
from descending to the orientation character and obstructs the
section-level identification
$\operatorname{pf}_X=\Delta_5$ on the wall complement.

\textit{Conditional theorem.}  Conditional on
\((O2_{\mathrm{atlas}})\) in Definition~\ref{def:G-O2-atlas},
Theorem~\ref{thm:G-O2-orbit-transport}.  The local rank-one crossing
model is established (Proposition~\ref{prop:rank-one-crossing-o2-sign});
component, boundary, and overlap compatibility on the full
half-Hilbert orbit is part of the retained wall-atlas datum and is
transported by Theorem~\ref{thm:G-O2-orbit-transport}.  The
$4096=2^{12}$ sign count is the scalar shadow of that retained datum,
as separated in Appendix~\ref{app:G-O2-discharge}.  The full
normal-form theorem with this sign count is relative to
Appendix~\ref{app:pfaffian-wall-normal-form}.

\subsubsection*{Joint theorem with recognition datum}

\begin{theorem}[Pfaffian--Dirac identification with primitive recognition datum]
\label{thm:pfaffian-dirac-conditional}
On \(K3\times E\), suppose (D0), the retained
\((O1_{\mathrm{quot}})\) datum, the \((O1)^+\) Weyl lifts, and
\((O2_{\mathrm{atlas}})\) hold at every finite HN height $R$ with
strict HN transitions as in
Appendix~\ref{app:obstruction-discharges}.  Suppose the finite
geometric Pfaffian datum \((P_{\mathrm{fin}})\) of
Definition~\ref{def:finite-geometric-pfaffian-datum} is supplied on a
cofinal HN system.  Suppose the recognition data (R0)--(R8) of
\S\ref{subsec:primitive-recognition} are supplied on a cofinal sequence
of saturated downward finite windows.  Suppose the chiral Koszul source
datum of \S\ref{subsec:koszul-source} is supplied, together with the
finite Koszul comparison datum of
Definition~\ref{def:finite-koszul-comparison-datum} and the strict
Koszul Mittag--Leffler exactness hypothesis.
Then
\[
\operatorname{Pf}_{\mathrm{prot}}(\mathfrak D_X^{\mathrm{DI}})
=\Delta_5,
\qquad
\varepsilon_o=\nu_{\Delta_5}|_{W^{(2)}(\Lambda_{II}^{2,1})},
\qquad
P_X\cong\mathfrak g_{\Delta_5},
\]
and $\Omega_E^{\mathrm{ch}}C_X\simeq\mathsf{Den}_{\Delta_5,E}$.
\end{theorem}

The eight finite constructions of
\S\ref{subsec:eight-finite-constructions}, the recognition theorem
of \S\ref{subsec:primitive-recognition}, and the Koszul comparison
of \S\ref{subsec:koszul-source} give the displayed implication.  The Pfaffian and orientation
conclusion holds by the retained D0-HN datum, the retained
\((O1_{\mathrm{quot}})\) datum, the chosen Weyl lifts, and the
retained \((O2_{\mathrm{atlas}})\) datum; the primitive
Lie-superalgebra and current-envelope conclusions require the cofinal
recognition datum, the Koszul source datum, and the finite Koszul
comparison datum.
The current-envelope conclusion uses the finite comparison
compatibility of
Proposition~\ref{prop:finite-koszul-comparison-compatibility}; without
that datum, the source cobar quasi-isomorphism need not preserve the
Weyl action, Pfaffian orientation character, Hall pairing, radical
quotient, or PBW filtration used in recognition.

\subsubsection*{The 4096 sign sum and the (O2) wall atlas}

The constant $4096=2^{12}$ is the size of the half-Hilbert
orbit of the type-II reflection group on the retained finite-window
components and boundary strata of the type-II wall configuration.
At the local rank-one chart, each retained component contributes a
sign $\varepsilon^{\mathrm{loc}}_{\delta_i,R}=-1$; the global sign
system is the retained (O2) wall atlas of Appendix~E transported along
the chosen Weyl lifts, and its character is
\(\nu_{\Delta_5}|_{W^{(2)}(\Lambda_{II}^{2,1})}\).
The Oberdieck--Pixton branch
$Z^X_{\mathrm{OP}}=-4096\,\Delta_5^{-2}$ is a scalar convention and is
not the source of the orientation character.  The equality
$\varepsilon_o=\nu_{\Delta_5}$ on
$W^{(2)}(\Lambda_{II}^{2,1})$ is relative to the retained quotient
orientation, the chosen Weyl lifts, and the (O2) wall atlas; the
local rank-one Pfaffian computation gives only the seed sign.

\subsubsection*{Compatibility with the companion volumes}

The four obstruction classes are source-side geometric data.  Their
automorphic shadow is smaller: the Gritsenko--Cl\'ery catalogue fixes
the Borcherds products, weights, characters, and denominator rows; it
does not record the \(K3\times E\) quotient orientation, the
Weyl-equivariant determinant-line lifts, or the Pfaffian wall atlas.
Thus (D0) is the cosection-reduced \(K3\)-fibre degeneration datum,
(O1) is the reduced orientation gerbe with finite-stabilizer descent,
(O1)\(^+\) is the Weyl-equivariant determinant-line transport, and
(O2) is the retained type-II Pfaffian wall atlas.  The
relation-closed window tests compatibility of these source data with
the target BKM relations; it is not an equality theorem between
geometric obstruction classes and automorphic catalogue rows.

The inverse limit on the HN system \(\mathcal R_{\mathrm{HN}}\)
contains the hybrid carrier, the cosection-twisted orientation, the
Hall correspondences, the protected Pfaffian, the Koszul source
coalgebra, the primitive recognition apparatus, the finite constructions,
and the retained orientation and wall-atlas data into
\[
\mathfrak D_X^{\mathrm{DI}}
=\varprojlim_R(\mathcal A_{X,R}^{E_3},
F_{X,R}^{\mathrm{hyb}},\dots,\mathrm{Rec}_R).
\]
The Pfaffian and orientation part of
\(\textcircled{2}\leftrightarrow\textcircled{3}\) is then the
retained-data theorem of Chapter~\ref{ch:pfaffian-dirac}; the
Lie-superalgebra recognition remains conditional on \((S1)\)--\((S10)\).
